\begin{center}
\LARGE
The rif and radical rif algorithms
\end{center}

\begin{center}
\Large
Gregory J. Reid
\end{center}

\noindent
Date:  July 17th (Wednesday)\\
Time:  10:30-11:00\\

\begin{center}
\large
Abstract
\end{center}

The rif algorithm uses a finite number of differentiations and algebraic operations to simplify
analytic nonlinear systems of partial differential equations to what we call reduced involutive
form. This form includes the integrability conditions of the system and satisfies a constant rank
condition. 

The rif algorithm is effective provided it can effectively perform elimination on the algebraic
systems it generates and transform such systems to constant rank form. The radical rif algorithm is
a realization of the rif algorithm for polynomially nonlinear pde systems which uses Buchberger's
algorithm for elimination and algorithms for constructing the radical of a polynomial ideal to
realize the constant rank condition. Various applications including the symmetry analysis of
differential equations are discussed. 

Reference:
G. J. Reid, A. D. Wittkopf and A. Boulton (1994).
Reduction of systems of nonlinear partial differential equations to simplified involutive forms. \\
Available at http://www.iam.ubc.ca/tr/1994/iam94-14. 

